\section{Conclusion}

This paper proposed a completion-oriented resource contract for the
serving layer, built it from four lessons packet networking paid for,
stated in advance and in public the numbers that would prove its
economics wrong, and then measured them. The economics are wrong: a
median 10.3 percent loss, no supporting cell in 108, the
graceful-degradation hypothesis reversed, and the cause structural, a
workload whose 5 percent duty cycle strands whatever is reserved for
it. What the contract provides regardless, isolation, containment, an
audit record, and a stop that the computation being stopped cannot
reach, is measured or holds by construction, and is now priced rather
than free. The starvation of compliant agentic work by request-oriented
schedulers is no longer an argument but a table. The evaluation leaves
one rule for the benchmarks being built around long-horizon agents: the
horizon must exceed the workload's completable lifetimes, or the
benchmark measures its own window. Every artifact required to check,
retune, or overturn these numbers is released.
